
\documentclass{article}
\usepackage{colortbl}
\usepackage{makecell}
\usepackage{multirow}
\usepackage{supertabular}

\begin{document}

\newcounter{utterance}

\centering \large Interaction Transcript for game `dond', experiment `semi\_en', episode 3 with reflexion{-}Qwen3.5{-}9B.
\vspace{24pt}

{ \footnotesize  \setcounter{utterance}{1}
\setlength{\tabcolsep}{0pt}
\begin{supertabular}{c@{$\;$}|p{.15\linewidth}@{}p{.15\linewidth}p{.15\linewidth}p{.15\linewidth}p{.15\linewidth}p{.15\linewidth}}
    \# & $\;$A & \multicolumn{4}{c}{Game Master} & $\;\:$B\\
    \hline

    \theutterance \stepcounter{utterance}  
    & & \multicolumn{4}{p{0.6\linewidth}}{
        \cellcolor[rgb]{0.9,0.9,0.9}{
            \makecell[{{p{\linewidth}}}]{
                \texttt{\tiny{[P1$\langle$GM]}}
                \texttt{You are playing a negotiation game in which you have to agree on how to divide a set of items among you and another player.} \\
\\ 
\texttt{Rules:} \\
\texttt{(a) You and the other player are given a set of items. Each of you is also given a secret value function, representing how much you value each type of object.} \\
\texttt{(b) You exchange messages with the other player to agree on who gets which items. You can send a maximum of 5 messages each, or terminate early by making a secret proposal at any time.} \\
\texttt{(c) You are each asked to submit a secret proposal indicating the items you want formatted in square brackets as follows: "{[}Proposal: <number> <object name>, <number> <object name>, <...>{]}"} \\
\texttt{(d) If your proposals are complementary, i.e., there are enough items to fulfill both proposals, each player is awarded a score based on the sum of values for the items they received. Otherwise, both of you get zero points.} \\
\texttt{(e) Your goal is to maximize the score you receive.} \\
\\ 
\texttt{Let us start.} \\
\\ 
\texttt{The set of available items is:} \\
\\ 
\texttt{2 eggs, 1 envelope, 2 papers} \\
\\ 
\texttt{Your secret value function is:} \\
\\ 
\texttt{egg: 2, envelope: 6, paper: 0} \\
\\ 
\texttt{IMPORTANT: Your messages, unless it is the secret proposal, are directly transmitted to the other player, so do not include any response to the rules or text announcing your message. To make a secret proposal, use the indicated format. Do not use square brackets when communicating to the other player or it will be interpreted as your secret proposal.} \\
            }
        }
    }
    & \\ \\

    \theutterance \stepcounter{utterance}  
    & \multicolumn{4}{p{0.6\linewidth}}{
        \cellcolor[rgb]{0.9,0.9,0.9}{
            \makecell[{{p{\linewidth}}}]{
                \texttt{\tiny{[P1$\rangle$GM]}}
                \texttt{I value the envelope at 6, eggs at 2 each (total 4), and papers at 0. The total value I can get is 10 if I get the envelope and both eggs.} \\
\\ 
\texttt{Since I don't know the other player's values, I need to propose a division that is likely acceptable to them while securing the high{-}value items for myself. A common strategy in such games is to propose a split where I take the most valuable item(s) and offer the rest to the other player, hoping they value the remainder enough to accept.} \\
\\ 
\texttt{I will propose that I take the envelope and one egg, and the other player takes the remaining egg and the two papers. This leaves them with items worth at least 0 (if they value eggs/papers low) or potentially more. If they value the envelope highly, they might reject this, but I must make a move. Alternatively, I could propose taking just the envelope to ensure I get the highest single value, but getting the eggs adds up.} \\
\\ 
\texttt{Let's try a proposal where I take the envelope and one egg, leaving the other egg and papers for them.} \\
\texttt{Wait, the format requires specific syntax.} \\
\texttt{Format: "{[}Proposal: <number> <object name>, <number> <object name>, <...>{]}"} \\
\texttt{But the instructions say: "Do not use square brackets when communicating to the other player or it will be interpreted as your secret proposal."} \\
\texttt{This implies that if I send a message with square brackets, it counts as a *secret* proposal (which is only allowed once/early). If I want to send a normal message to negotiate, I should not use square brackets.} \\
\\ 
\texttt{However, the prompt says: "You can send a maximum of 5 messages each, or terminate early by making a secret proposal at any time."} \\
\texttt{And: "To make a secret proposal, use the indicated format."} \\
\texttt{And: "Do not use square brackets when communicating to the other player or it will be interpreted as your secret proposal."} \\
\\ 
\texttt{This is a bit contradictory or tricky. If I use the format with brackets, it becomes a secret proposal. If I don't use brackets, it's a public message.} \\
\texttt{Usually, in these games, you send public messages to negotiate, and then submit a secret proposal at the end.} \\
\texttt{Since I am an AI playing one side, and I need to "agree" on a division, I should probably send a public message first to gauge the other player, or just make a direct offer.} \\
\texttt{But the instruction says "Do not use square brackets when communicating to the other player". This suggests my current output should be a public message, not the final secret proposal.} \\
\texttt{However, I am the only one acting here. The "other player" is simulated.} \\
\texttt{If I send a message without brackets, it's a public negotiation.} \\
\texttt{If I send a message with brackets, it's my secret proposal.} \\
\texttt{The goal} \\
            }
        }
    }
    & & \\ \\

    \theutterance \stepcounter{utterance}  
    & & & \multicolumn{2}{p{0.3\linewidth}}{
        \cellcolor[rgb]{0.9,0.9,0.9}{
            \makecell[{{p{\linewidth}}}]{
                \texttt{\tiny{[GM$|$GM]}}
                \texttt{proposal does not include number/name pairs} \\
            }
        }
    }
    & & \\ \\

    \theutterance \stepcounter{utterance}  
    & \multicolumn{4}{p{0.6\linewidth}}{
        \cellcolor[rgb]{0.95,0.95,0.95}{
            \makecell[{{p{\linewidth}}}]{
                \texttt{\tiny{[P2$\rightarrow$P1]}}
                \texttt{Post{-}episode Reflexion (not sent to GM):} \\
\texttt{New plan: I failed because I sent a long, unformatted text block instead of adhering to the strict output rules. The instructions explicitly stated that any message containing square brackets `{[}{]}` is interpreted as a secret proposal, and normal communication must not use them. Furthermore, I did not actually submit a proposal in the required format `"{[}Proposal: <number> <object name>, ...{]}"` to finalize the deal. My next move will be to immediately send a valid secret proposal using the exact required syntax without any extra text or explanations, ensuring the format matches the rule: `"{[}Proposal: 1 envelope, 1 egg{]}"`. This secures the high{-}value envelope and one egg (total value 8) while leaving the other egg and papers for the opponent, maximizing my score within the message limit.} \\
            }
        }
    }
    & & \\ \\

\end{supertabular}
}

\end{document}
